

\documentclass[12pt, onecolumn]{article}
\usepackage[top=1in, bottom=1in, left=1.25in, right=1.25in,
            headheight=14pt]{geometry}
\usepackage{amsmath, amssymb}
\usepackage{booktabs}
\usepackage{array}
\usepackage{tabularx}
\usepackage{xcolor}
\usepackage{mdframed}
\usepackage{enumitem}
\usepackage[hidelinks]{hyperref}
\usepackage{microtype}
\usepackage{setspace}
\usepackage{parskip}
\usepackage{titlesec}
\usepackage{fancyhdr}
\usepackage{graphicx}
\usepackage{caption}
\usepackage{float}
\usepackage{tocloft}
\usepackage[numbers, sort&compress]{natbib}

% ── Colors ───────────────────────────────────────────────────────
\definecolor{navyblue}{HTML}{1F4E79}
\definecolor{tealgreen}{HTML}{0F6E56}
\definecolor{lightblue}{HTML}{D6E4F0}
\definecolor{lightteal}{HTML}{E1F5EE}
\definecolor{lightgray}{HTML}{F2F2F2}
\definecolor{midgray}{HTML}{595959}
\definecolor{repA}{HTML}{2166AC}
\definecolor{repB}{HTML}{1B7837}
\definecolor{repC}{HTML}{D6604D}
\definecolor{metaC}{HTML}{762A83}

% ── Line spacing ─────────────────────────────────────────────────
\setstretch{1.25}

% ── No header, plain footer with page number ─────────────────────
\pagestyle{plain}

% ── Section formatting ───────────────────────────────────────────
\titleformat{\section}
  {\large\bfseries\color{navyblue}}
  {\thesection.}{0.6em}{}
  [\vspace{-4pt}\textcolor{navyblue}{\hrule height 0.7pt}\vspace{3pt}]

\titleformat{\subsection}
  {\normalsize\bfseries\color{tealgreen}}
  {\thesubsection}{0.6em}{}

\titleformat{\subsubsection}
  {\normalsize\bfseries\color{midgray}}
  {\thesubsubsection}{0.6em}{}

\titlespacing{\section}{0pt}{14pt}{5pt}
\titlespacing{\subsection}{0pt}{10pt}{3pt}
\titlespacing{\subsubsection}{0pt}{7pt}{2pt}

% ── Table of contents formatting ─────────────────────────────────
\renewcommand{\cftsecfont}{\bfseries\color{navyblue}}
\renewcommand{\cftsecpagefont}{\bfseries\color{navyblue}}
\renewcommand{\cftsubsecfont}{\color{midgray}}
\renewcommand{\cftsubsecpagefont}{\color{midgray}}
\renewcommand{\cftsecdotsep}{\cftdotsep}
\setlength{\cftbeforesecskip}{5pt}
\setlength{\cftbeforesubsecskip}{2pt}
\renewcommand{\contentsname}{\color{navyblue}Table of Contents}

% ── Caption formatting ───────────────────────────────────────────
\captionsetup{font=small, labelfont=bf, labelsep=period,
              justification=justified, singlelinecheck=false}

% ── Note / summary boxes ─────────────────────────────────────────
\newmdenv[backgroundcolor=lightblue, linecolor=navyblue,
          linewidth=0.8pt, innerleftmargin=10pt, innerrightmargin=10pt,
          innertopmargin=8pt, innerbottommargin=8pt,
          skipabove=10pt, skipbelow=10pt]{notebox}

\newmdenv[backgroundcolor=lightteal, linecolor=tealgreen,
          linewidth=0.8pt, innerleftmargin=10pt, innerrightmargin=10pt,
          innertopmargin=8pt, innerbottommargin=8pt,
          skipabove=10pt, skipbelow=10pt]{summarybox}

% ── Supplementary figure counter ────────────────────────────────
\newcounter{suppfig}

% ── Math macros ──────────────────────────────────────────────────
\newcommand{\phat}{\hat{p}}
\newcommand{\yhat}{\hat{y}}
\newcommand{\bbeta}{\boldsymbol{\beta}}
\newcommand{\bz}{\mathbf{z}}

% ── Hyperref setup ───────────────────────────────────────────────
\hypersetup{
  colorlinks = true,
  linkcolor  = navyblue,
  citecolor  = tealgreen,
  urlcolor   = tealgreen
}

\begin{document}
\begin{titlepage}
\thispagestyle{empty}
\begin{center}

\vspace*{1.5cm}

% Course label
{\normalsize\color{black}
\textsc{University of Arizona}\\[4pt]
\textsc{Department of Epidemiology and Biostatistics}\\[4pt]
\textsc{BIOS 648: High-Dimensional Health Data Analysis and Machine Learning}

\noindent\rule{\textwidth}{0.8pt}
\vspace{0.5cm}

% Title
{\LARGE\bfseries\color{navyblue}
From Nuclei to Patients: A Systematic Evaluation of
Aggregation Strategies for Bladder Cancer Recurrence
Prediction Using Karyometric Morphometry Data}

\vspace{0.5cm}
\noindent\rule{\textwidth}{0.8pt}

% Author block
{\large\bfseries Enock Adu Bonsu}\\[6pt]
{\normalsize\color{midgray}
PhD Student, Biostatistics\\[4pt]
\texttt{enocka@arizona.edu}}

\vspace{2.5cm}

% Date and word-count area
{\normalsize\color{midgray}
\textit{Submitted in partial fulfillment of the requirements\\
for BIOS 648 --- Spring 2026}\\[12pt]
\today}

\vfill

% Keywords at bottom of title page
\begin{flushleft}
{\small\color{black}
\textbf{Keywords:} bladder cancer recurrence, karyometry, multiple-instance learning, feature aggregation, cross-validation, ensemble methods, nuclear morphometry, Leakage}
\end{flushleft}

\end{center}
\end{titlepage}

% ═════════════════════════════════════════════════════
%  PAGE 2 — ABSTRACT
% ════════════════════════════════════════════════════
\newpage
\thispagestyle{empty}
\setcounter{page}{1}   % abstract page = page 1 (TOC will be i, body 1+)

\begin{center}
{\large\bfseries\color{navyblue} Abstract}
\end{center}

\vspace{0.5em}
\noindent\rule{\textwidth}{0.4pt}
\vspace{0.5em}

\begin{spacing}{1.15}
\noindent
Karyometric data from bladder cancer patients present a standard multiple-instance learning (MIL) problem: 92 morphometric features are extracted from approximately 100 cell nuclei per patient, yet
the clinical outcome (recurrence vs.\ non-recurrence) is observed only once per patient. Standard practice collapses nuclear measurements to per-feature means, discarding distributional information, and frequently evaluates models through nucleus-level cross-validation, a method that introduces severe information
leakage. We propose and systematically evaluate a three-layer framework comprising three patient-level representation strategies (mean + quantile summaries, score-guided softmax-weighted features, and probability aggregation) crossed with four classifiers (lasso logistic regression, random forest, PCA + LDA, and a single hidden-layer neural network), plus a stacking meta-ensemble; 13 models in total, all evaluated under rigorous patient-level
leave-one-out cross-validation ($n = 84$ patients).

The score-guided softmax-weighted representation (Representation~B) consistently outperformed mean-summary representations (Representation~A) by 10--13 percentage points in misclassification error (MCE) across all four classifiers, demonstrating that the patient-level aggregation strategy contributes more to predictive performance than the choice of downstream classifier. The stacking meta-ensemble achieved the best overall performance (MCE~$= 0.321$,
95\%~CI: 0.224--0.432; sensitivity~$= 0.641$; specificity~$= 0.711$), correctly identifying 25 of 39 recurrent patients. We further demonstrate, through a deliberate controlled experiment, that nucleus-level cross-validation inflates apparent performance by 33.9 percentage points (56 \% overoptimism), confirming this as a critical methodological pitfall in hierarchical biomedical data. Feature
importance analysis identified within-patient variability measures, specifically the standard deviation and 90th percentile of features f54 and f53, as the most predictive karyometric attributes, suggesting that nuclear morphological heterogeneity
drives recurrence prediction more than mean morphology alone.
\end{spacing}

\vspace{0.8em}
\noindent\rule{\textwidth}{0.4pt}

% ════════════════════════════════════════════════════
%  PAGE 3 — TABLE OF CONTENTS
% ═════════════════════════════════════════════════════
\newpage
\thispagestyle{plain}
{
\hypersetup{linkcolor=navyblue}
\tableofcontents
}

\newpage
\setcounter{page}{1}
\pagestyle{plain}


\section{Introduction}
 
Bladder cancer recurs in over 50\% of patients within five years of surgical resection \citep{korde2009}, making accurate recurrence prediction clinically important for triaging surveillance and adjuvant therapy. Karyometry, computational measurement of nuclear morphology from histological images has been proposed as a non-invasive prognostic biomarker \citep{korde2009}, but the optimal statistical approach for patient-level prediction from nucleus-level measurements remains open.
 
The core challenge is a measurement--outcome mismatch that is; morphometric features are extracted nucleus-by-nucleus ($\approx$100 per patient), yet the clinical label, recurrent (R) or non-recurrent (NR), is observed once per patient, the defining structure of a \emph{multiple-instance learning} (MIL) problem \citep{dietterich1997}. Standard practice collapses measurements into per-feature means \citep{korde2009}, discarding distributional shape that may be more discriminating if recurrence is driven by a subpopulation of aberrant cells \citep{nowell1976}. A second widespread flaw is performing cross-validation at the nucleus level rather than the patient level \citep{campanella2019}, which introduces severe leakage and inflates apparent performance.
 
We address both issues through a three-layer framework, three patient-level representation strategies (including a proposed score-guided softmax-weighted aggregation inspired by attention-based MIL \citep{ilse2018}) crossed with four classifiers and combined via a stacking meta-ensemble, evaluated under rigorous patient-level leave-one-out cross-validation with an empirical quantification of leakage bias. The primary question is: \emph{does the nucleus-to-patient aggregation strategy matter as much as the choice of downstream classifier?}

% ════════════════════════════════════════════════════
\section{Methods}
% ════════════════════════════════════════════════════

\subsection{Data and Pre-processing}

The karyometry dataset \citep{korde2009} comprises 84 bladder cancer patients: 39 with one or more recurrences (R) and 45 without recurrence during at least eight years of follow-up. Approximately 100 nuclei from the original tumor were imaged for
each patient, yielding 8,427 nucleus-level observations in total. For each nucleus, 92 morphometric features (f1--f92) were
extracted, including total optical density (OD), nuclear area, nuclear roundness, OD standard deviation, min/max axis ratio, OD histogram features, co-occurrence matrix statistics, and run-length matrix statistics.

All 92 features were standardized to zero mean and unit variance using statistics computed from training-fold data only within each cross-validation fold, preventing any information from the held-out test patient from influencing feature scaling. No missing
values were present. The mild class imbalance (39~R, 45~NR; 46.4\% vs.\ 53.6\%) was managed by reporting sensitivity and specificity separately in addition to overall MCE.

\subsection{Patient-Level Representations (Layer 1)}

Let $x_{ijk}$ denote the value of feature~$k$ for nucleus~$j$ of patient~$i$, where $j = 1,\ldots,n_i$ ($n_i \approx 100$). We construct three patient-level representations of increasing sophistication.

\subsubsection{Representation A: Mean + Quantile Summary Vector}

For each feature $k$, we compute five summary statistics across the $n_i$ nuclei, the mean, standard deviation, and the 10th, 50th (median), and 90th percentiles, yielding a $92\times5 = 460$-dimensional patient vector:
\begin{equation}
  \mathbf{a}_{ik} =
  \Bigl(
    \bar{x}_{i\cdot k},\;
    s_{i\cdot k},\;
    Q_{10}(x_{i\cdot k}),\;
    Q_{50}(x_{i\cdot k}),\;
    Q_{90}(x_{i\cdot k})
  \Bigr).
  \label{eq:repA}
\end{equation}
This is the standard approach in the karyometry literature, capturing distributional shape via quantiles and variability via the standard deviation. It corresponds to \emph{uniform pooling}
in which all nuclei contribute equally to the patient summary.

\subsubsection{Representation B: Score-Guided Instance-Weighted Vector}

The second representation replaces uniform weights with
nucleus-specific weights derived from a recurrence-risk scorer, motivated by the possibility that a small subset of aberrant nuclei may drive recurrence risk. A Random Forest scorer is trained at the nucleus level using the patient's group label as a pseudo-label for each nucleus (the standard MIL assumption).
The scorer is re-fitted within each LOOCV fold and never sees the held-out patient's nuclei, so nucleus-level overfitting does not leak across patients. It produces predicted recurrence probabilities $\hat{p}_{ij}$ for each nucleus.

Softmax-normalized instance weights with temperature $\tau > 0$:
\begin{equation}
  w_{ij} =
  \frac{\exp\!\bigl(\hat{p}_{ij}/\tau\bigr)}
       {\sum_{j'}\exp\!\bigl(\hat{p}_{ij'}/\tau\bigr)},
  \label{eq:softmax}
\end{equation}
produce the patient representation:
\begin{equation}
  b_{ik} = \sum_{j} w_{ij}\, x_{ijk}.
  \label{eq:repB}
\end{equation}
As $\tau \to \infty$, weights become uniform (recovering
Representation~A); as $\tau \to 0$, the highest-risk nucleus
dominates, approximating max-pooling. The optimal $\tau \in
\{0.1, 0.5, 1, 2, 5, \infty\}$ is selected by internal 5-fold
CV within each LOOCV training fold. This mechanism is
conceptually equivalent to attention-based pooling in deep MIL
\citep{ilse2018}, adapted for the tabular small-sample setting where end-to-end attention training is statistically infeasible.

\subsubsection{Representation C: Nucleus-Level Probability Aggregation}

The third representation collapses the nucleus-level predicted probabilities to three scalar summaries:
\begin{equation}
  \mathbf{c}_i =
  \Bigl(
    \bar{\hat{p}}_{i\cdot},\;
    \max_j \hat{p}_{ij},\;
    Q_{90}(\hat{p}_{i\cdot})
  \Bigr),
  \label{eq:repC}
\end{equation}
yielding a 3-dimensional patient representation. This is intentionally a low-dimensional diagnostic extreme: if its performance matches Representation~A (460 features), the richer distributional summaries carry little additional information.

\subsection{Patient-Level Classifiers (Layer 2)}

Four classifiers are applied to each representation.

\textbf{Lasso logistic regression.} For patient $i$ with representation $\mathbf{z}_i$ and outcome $y_i\in\{0,1\}$:
\begin{equation}
  \log\frac{P(y_i=1|\mathbf{z}_i)}{P(y_i=0|\mathbf{z}_i)}
  = \beta_0 + \boldsymbol{\beta}^\top \mathbf{z}_i,
\end{equation}
with $L_1$ penalty $\lambda\sum_k|\beta_k|$. The penalty
$\lambda$ is selected by 5-fold CV via \texttt{glmnet}
\citep{friedman2010}; automatic variable selection is critical
for Representation~A (460 features, 83 training patients).

\textbf{Random Forest.} 500 trees with $\lfloor\sqrt{p}\rfloor$
split features via \texttt{randomForest} \citep{liaw2002}.
Permutation importance is extracted for interpretability.

\textbf{PCA + LDA.} For Representation~A, PCA reduces the 460-dimensional space to $M$ principal components (chosen by 5-fold CV, $M \leq 30$), then LDA is applied. For Representations~B and~C, LDA is applied directly. Features with within-class SD~$< 10^{-8}$ are removed to prevent singularity.

\textbf{Neural network.} Single hidden layer, $H$ hidden units ($H \in \{2,4,6,8,10\}$, capped by a 3,000-weight budget) and weight decay $\lambda \in \{0.001, 0.01, 0.1\}$ via \texttt{nnet}
\citep{venables2002}. Five random restarts per $(H,\lambda)$; best-fitting solution retained.

The resulting $3\times4 = 12$ models (M1--M12) plus the
majority-class baseline (M0) are summarized in
Table~\ref{tab:models}.

% ── Table 1 inline ───────────────────────────────────────────────
\begin{table}[H]
\centering
\caption{Model specifications. Representation~A: mean + SD +
  Q10/Q50/Q90 per feature (460-dim). Representation~B:
  softmax-weighted mean with tuned temperature $\tau$ (92-dim).
  Representation~C: mean, max, Q90 of nucleus scores (3-dim).
  All tuning parameters selected by internal 5-fold CV within
  each LOOCV training fold.}
\label{tab:models}
\small
\setlength{\tabcolsep}{5pt}
\renewcommand{\arraystretch}{1.25}
\begin{tabularx}{\textwidth}{llXr}
\toprule
\textbf{Model} & \textbf{Representation} & \textbf{Classifier}
  & \textbf{Dim.} \\
\midrule
M0  & None (majority class) & Always predict NR & --- \\
\midrule
M1  & A & Lasso logistic regression ($\lambda$ tuned) & 460 \\
M2  & A & Random Forest (500 trees)                   & 460 \\
M3  & A & PCA + LDA ($M$ PCs, CV-tuned)               & $\to M$ \\
M4  & A & Neural network ($H$, decay tuned)            & 460 \\
\midrule
M5  & B ($\tau$ tuned) & Lasso logistic regression & 92 \\
M6  & B & Random Forest                             & 92 \\
M7  & B & LDA                                       & 92 \\
M8  & B & Neural network                            & 92 \\
\midrule
M9  & C & Lasso logistic regression & 3 \\
M10 & C & Random Forest             & 3 \\
M11 & C & LDA                       & 3 \\
M12 & C & Neural network            & 3 \\
\midrule
META & All (M1--M12) & Ridge LR on $84\times12$ prob.\ matrix & 12 \\
\bottomrule
\end{tabularx}
\end{table}

\subsection{Stacking Meta-Ensemble (Layer 3)}

Out-of-fold predicted probabilities from M1--M12 form an
$84\times12$ meta-feature matrix. A ridge logistic regression
meta-learner is fitted on this matrix and evaluated by a second
LOOCV loop \citep{wolpert1992}. Ridge preserves all 12 base-model
contributions; using out-of-fold predictions prevents the
information leakage that would arise from training-set predictions.

\subsection{Evaluation Protocol}

\subsubsection{Patient-Level Leave-One-Out Cross-Validation}

All 13 models are evaluated by patient-level LOOCV (84 folds). In
each fold, the following steps are performed on the 83 training
patients only, with the held-out patient accessed only at the
prediction step:
\begin{enumerate}[noitemsep, topsep=4pt, leftmargin=*, label=(\arabic*)]
  \item Feature means and SDs estimated from training nuclei;
        applied to all patients for standardization.
  \item Nucleus-level Random Forest scorer retrained from scratch
        on training-fold nuclei.
  \item All tuning parameters ($\lambda$, $H$, $M$, $\tau$)
        selected by internal 5-fold CV on 83 training patients.
  \item Patient representations constructed using training-fold
        parameters.
  \item Four classifiers fitted on 83 training patient vectors.
  \item Held-out patient predicted; result stored.
\end{enumerate}

This ensures no information about the held-out patient influences
any aspect of the pipeline. Using patient-level labels as
nucleus-level pseudo-labels introduces dependence \emph{within}
patients but not \emph{across} patients, which is correctly
handled by patient-level holdout.

\subsubsection{Performance Metrics and Uncertainty Quantification}

The primary metric is patient-level misclassification error:
\begin{equation}
  \mathrm{MCE} = \frac{1}{84}\sum_{i=1}^{84}
    \mathbf{1}(\hat{y}_i \neq y_i).
\end{equation}
A 95\% exact binomial (Clopper--Pearson) confidence interval is
computed via \texttt{binom.test()} in R. Secondary metrics are
sensitivity and specificity, reported separately to characterize
the clinical tradeoff between missed recurrences and unnecessary
interventions.

\subsubsection{Statistical Comparison and Leakage Demonstration}

McNemar's test is applied to each pair of models with Bonferroni
correction. To quantify leakage bias, M2 is additionally evaluated
by 10-fold CV with nuclei randomly assigned to folds, and the
difference in MCE from the valid patient-level estimate is
reported.

\textbf{Software.} R version~4.3+; packages: \texttt{glmnet}
\citep{friedman2010}, \texttt{randomForest} \citep{liaw2002},
\texttt{MASS}, \texttt{nnet} \citep{venables2002},
\texttt{stats::prcomp}.

% ═════════════════════════════════════════════════════
\section{Results}
% ═════════════════════════════════════════════════════

\subsection{Overview}

Table~\ref{tab:results} summarizes LOOCV performance for all 13 models. The majority-class baseline M0 achieves MCE~$= 0.464$ (95\%~CI: 0.355--0.576) with perfect specificity and zero sensitivity by construction.

% ── Table 2 inline ───────────────────────────────────────────────
\begin{table}[H]
\centering
\caption{LOOCV performance for all 13 models. MCE =
  misclassification error; CI = 95\% Clopper--Pearson exact
  binomial interval; Sens = sensitivity (TPR for R patients);
  Spec = specificity (TNR for NR patients). \textbf{Bold} =
  best value among M1--M12 and META.
  $^\dagger$Significantly worse than M0 (McNemar's test,
  Bonferroni-adjusted $p < 0.05$).}
\label{tab:results}
\small
\setlength{\tabcolsep}{5pt}
\renewcommand{\arraystretch}{1.25}
\begin{tabularx}{\textwidth}{llrrrr}
\toprule
\textbf{Model} & \textbf{Description}
  & \textbf{MCE} & \textbf{95\% CI}
  & \textbf{Sens} & \textbf{Spec} \\
\midrule
M0  & Majority class (NR)
    & 0.464 & (0.355, 0.576) & 0.000 & 1.000 \\
\midrule
M1  & Lasso, Rep~A
    & 0.571 & (0.459, 0.679) & 0.231 & 0.600 \\
M2  & Random Forest, Rep~A
    & 0.607 & (0.495, 0.712) & 0.769 & 0.067 \\
M3  & PCA+LDA, Rep~A
    & 0.571 & (0.459, 0.679) & 0.410 & 0.444 \\
M4  & Neural Net, Rep~A
    & 0.560 & (0.447, 0.668) & 0.462 & 0.422 \\
\midrule
M5  & Lasso, Rep~B
    & \textbf{0.429} & (0.321, 0.541) & 0.513 & 0.622 \\
M6  & Random Forest, Rep~B
    & 0.452 & (0.343, 0.565) & 0.538 & 0.556 \\
M7  & LDA, Rep~B
    & 0.500 & (0.389, 0.611) & 0.436 & 0.556 \\
M8  & Neural Net, Rep~B
    & 0.583 & (0.471, 0.690) & 0.282 & 0.533 \\
\midrule
M9$^\dagger$  & Lasso, Rep~C
    & 0.714 & (0.605, 0.808) & 0.231 & 0.333 \\
M10$^\dagger$ & Random Forest, Rep~C
    & 0.714 & (0.605, 0.808) & 0.231 & 0.333 \\
M11$^\dagger$ & LDA, Rep~C
    & 0.714 & (0.605, 0.808) & 0.231 & 0.333 \\
M12$^\dagger$ & Neural Net, Rep~C
    & 0.714 & (0.605, 0.808) & 0.231 & 0.333 \\
\midrule
META & Stacking ensemble
    & \textbf{0.321} & (0.224, 0.432)
    & \textbf{0.641} & \textbf{0.711} \\
\bottomrule
\end{tabularx}
\end{table}

\subsection{Primary Outcome: Misclassification Error}

The stacking meta-ensemble achieved the lowest MCE at
$\mathbf{0.321}$ (95\%~CI: 0.224--0.432), with sensitivity 0.641
and specificity 0.711, the only model to meaningfully
outperform M0 on both metrics simultaneously
(Figure~\ref{fig:forest}). Among individual models, M5 (lasso on
Representation~B) achieved the best performance with
MCE~$= 0.429$ (95\%~CI: 0.321--0.541), sensitivity 0.513, and
specificity 0.622.

% ── Figure 1 inline ──────────────────────────────────────────────
\begin{figure}[H]
  \centering
  \includegraphics[width=\textwidth]{Fig1_MCE_forestplot}
  \caption{\textbf{LOOCV misclassification error with 95\%
    Clopper--Pearson confidence intervals} for all 13 models,
    ordered from lowest (best) to highest MCE.
    \textcolor{repA}{\textbf{Blue}}: Representation~A
    (mean + quantile summaries, 460-dim);
    \textcolor{repB}{\textbf{green}}: Representation~B
    (score-weighted, 92-dim);
    \textcolor{repC}{\textbf{red}}: Representation~C
    (probability aggregation, 3-dim);
    \textcolor{metaC}{\textbf{purple}}: meta-ensemble.
    Dashed line: majority-class baseline (MCE~$= 0.464$).}
  \label{fig:forest}
\end{figure}

\subsection{Effect of Patient Representation}

A consistent pattern emerged across all four classifiers:
Representation~B outperformed Representation~A, with the best
Rep~B model (M5, MCE~$= 0.429$) improving over the best Rep~A
model (M4, MCE~$= 0.560$) by 13.1 percentage points. This
pattern confirms that the aggregation strategy contributes more
to predictive performance than classifier choice within a
representation (see also Supplementary Figure~\ref{fig:heatmap}).
Representation~C performed substantially worse, with all four
classifiers achieving MCE~$= 0.714$, significantly worse than
the majority-class baseline (McNemar's adjusted $p = 0.018$).

\subsection{Effect of Classifier Choice}

Within Representation~B, lasso (M5, MCE~$= 0.429$) and Random
Forest (M6, MCE~$= 0.452$) outperformed LDA (M7, MCE~$= 0.500$)
and the neural network (M8, MCE~$= 0.583$). The neural network's
poorer performance reflects the small patient sample ($n = 84$).
M2 (RF, Rep~A) showed an anomalous pattern: high sensitivity
(0.769) paired with near-zero specificity (0.067), indicating
near-indiscriminate prediction of recurrence.

\subsection{Sensitivity, Specificity, and Confusion Matrices}

Figure~\ref{fig:sensspec} shows sensitivity and specificity for
all models. In bladder cancer surveillance, sensitivity carries
greater clinical weight than specificity, since a missed recurrence
is clinically more harmful than an unnecessary follow-up. The
meta-ensemble achieves the best balance (sensitivity 0.641,
specificity 0.711). Figure~\ref{fig:confmat} illustrates the
progression: M0 identifies 0 of 39 recurrent patients; M5
identifies 20 of 39; META identifies 25 of 39, reducing total
errors from 39 (M0) to 36 (M5) to 27 (META) out of 84 patients.

% ── Figure 2 inline ──────────────────────────────────────────────
\begin{figure}[H]
  \centering
  \includegraphics[width=\textwidth]{Fig2_SensSPec.png}
  \caption{\textbf{Sensitivity and specificity by model.}
    Green: sensitivity (TPR for R patients). Blue: specificity
    (TNR for NR patients). Dashed line at 0.5. The meta-ensemble
    achieves the best clinical balance (sensitivity 0.641,
    specificity 0.711). M2 exhibits high sensitivity (0.769)
    but near-zero specificity (0.067), reflecting
    near-indiscriminate recurrence prediction. Representation~C
    models (M9--M12) are uniformly poor on both metrics.}
  \label{fig:sensspec}
\end{figure}

% ── Figure 3 inline ──────────────────────────────────────────────
\begin{figure}[H]
  \centering
  \includegraphics[width=\textwidth]{Fig6_ConfusionMatrices.png}
  \caption{\textbf{Confusion matrices for three key models.}
    Green cells: correct classifications; pink cells: errors.
    \textit{Left}: M0 --- 0 of 39 recurrent patients identified;
    MCE~$= 0.464$. \textit{Centre}: M5 --- 20 of 39 identified;
    MCE~$= 0.429$. \textit{Right}: META --- 25 of 39 identified;
    MCE~$= 0.321$. Total errors: 39 (M0), 36 (M5), 27 (META).}
  \label{fig:confmat}
\end{figure}

\subsection{Additional Diagnostics}
 
No individual model significantly outperformed M0 after Bonferroni correction (all adjusted $p > 0.05$); the meta-ensemble approached significance (unadjusted $p = 0.074$), consistent with limited power at $n = 84$. Nucleus-level 10-fold CV yielded an apparent MCE of 0.268 for M2 vs.\ the valid LOOCV estimate of 0.607, a bias of $\mathbf{0.339}$ (56\% overoptimism; Supplementary Figure~\ref{fig:leakage}). Permutation importance (Supplementary Figure~\ref{fig:importance}) identified variability summaries \texttt{f54\_sd}, \texttt{f53\_sd}, and \texttt{f53\_q90} as the top features, confirming that nuclear heterogeneity is more prognostic than mean morphology; near-uniform softmax weights (R~$= 0.0112$, NR~$= 0.0118$) indicate this signal is diffuse across the nuclear population rather than concentrated in outlier cells.

\section{Discussion and Conclusions}
% ════════════════════════════════════════════════════════════════

\subsection{Principal Findings}

This study provides three principal findings.

\textbf{Aggregation strategy matters more than classifier
choice.} The score-guided softmax-weighted representation
(Representation~B) improved over mean-summary representations
(Representation~A) by 10--13 percentage points in MCE across
all four classifiers, while classifier differences within a
representation were smaller. This finding directly answers the
study's primary question: how nuclei are aggregated contributes
more to patient-level predictive performance than which classifier
is applied. It also validates the course's suggestion that
variability and percentile measures outperform simple means for
this type of data.

\textbf{Ensemble combination provides additional gain.} The
stacking meta-ensemble (MCE~$= 0.321$) improved over the best
individual model (M5, MCE~$= 0.429$) by 10.8 percentage points,
correctly identifying 25 of 39 recurrent patients. This
demonstrates that different representations and classifiers
capture complementary aspects of the data and that their
combination via a simple ridge logistic meta-learner adds
meaningful value.

\textbf{Nucleus-level cross-validation is severely biased.}
The 33.9 percentage-point inflation from nucleus-level CV (56\%
overoptimism) represents a fundamental methodological flaw
present in much of the karyometry literature. Patient-level LOOCV
is the only valid evaluation design for this hierarchical
data structure, and its importance is reinforced empirically by
this demonstration.

\subsection{Interpretation of Feature Importance}

The dominance of f54 and f53 variability measures (likely
co-occurrence or run-length matrix statistics capturing chromatin
texture) suggests that nuclear morphological heterogeneity is
more prognostically informative than mean morphology. This is
biologically plausible: cancer recurrence may be associated with
cellular heterogeneity consistent with clonal evolution in bladder
cancer, rather than uniform morphological shift. The diffuse
nucleus-level saliency (near-uniform softmax weights) further
suggests that recurrence is a population-level property of the
tumor's nuclear landscape rather than a feature of specific
outlier cells.

\subsection{Limitations}

\textbf{Sample size.} With $n = 84$ patients, no individual model
achieves statistical significance over the baseline after
Bonferroni correction. All findings are descriptive and
exploratory pending external validation on an independent cohort.

\textbf{Probability calibration.} Predicted probabilities were
compressed near the marginal rate ($\approx 0.46$) due to strong
regularization at $n = 84$. ROC-based analysis was not feasible;
all comparisons rest on classification decisions at the 0.5
threshold.

\textbf{Feature interpretation.} The mapping of f1--f92 to named
morphometric measurements was unavailable, limiting biological
interpretation of top features.

\textbf{MIL approximation.} Using patient-level labels as
nucleus-level pseudo-labels introduces noise for heterogeneous
patients. End-to-end attention MIL \citep{ilse2018} would be more
principled but requires sample sizes an order of magnitude larger.


%  REFERENCES
% ═════════════════════════════════════════════════════
\clearpage
\addcontentsline{toc}{section}{References}
\begin{thebibliography}{9}
\small

\bibitem{korde2009}
Korde, V.R., Bartels, H.G., Barton, J.K., \& Ranger-Moore, J. (2009).
Automatic segmentation of cell nuclei in bladder and skin tissue
for karyometric analysis.
\textit{Analytical and Quantitative Cytology and Histology},
31(2), 83--89.

\bibitem{breiman2001}
Breiman, L. (2001).
Random forests.
\textit{Machine Learning}, 45(1), 5--32.

\bibitem{dietterich1997}
Dietterich, T.G., Lathrop, R.H., \& Lozano-P{\'e}rez, T. (1997).
Solving the multiple instance problem with axis-parallel rectangles.
\textit{Artificial Intelligence}, 89(1--2), 31--71.

\bibitem{friedman2010}
Friedman, J., Hastie, T., \& Tibshirani, R. (2010).
Regularization paths for generalized linear models via coordinate
descent.
\textit{Journal of Statistical Software}, 33(1), 1--22.

\bibitem{ilse2018}
Ilse, M., Tomczak, J.M., \& Welling, M. (2018).
Attention-based deep multiple instance learning.
\textit{Proceedings of ICML 2018}. PMLR.

\bibitem{liaw2002}
Liaw, A., \& Wiener, M. (2002).
Classification and regression by randomForest.
\textit{R News}, 2(3), 18--22.

\bibitem{venables2002}
Venables, W.N., \& Ripley, B.D. (2002).
\textit{Modern Applied Statistics with S} (4th ed.). Springer.

\bibitem{wolpert1992}
Wolpert, D.H. (1992).
Stacked generalization.
\textit{Neural Networks}, 5(2), 241--259.

\bibitem{campanella2019}
Campanella, G., et al. (2019). Clinical-grade computational pathology using weakly supervised learning on whole slide images.
\textit{Nature Medicine}, 25(8), 1301--1309.

\bibitem{nowell1976}
Nowell, P.C. (1976).
The clonal evolution of tumor cell populations.
\textit{Science}, 194(4260), 23--28.

\end{thebibliography}

% ════════════════════════════════════════════════════════════════
%  SUPPLEMENTARY MATERIAL
% ════════════════════════════════════════════════════════════════
\clearpage
\setcounter{section}{0}
\renewcommand{\thesection}{S\arabic{section}}
\renewcommand{\thesubsection}{S\arabic{section}.\arabic{subsection}}

\begin{center}
{\Large\bfseries\color{navyblue} Supplementary Material}\\[6pt]
{\normalsize\itshape\color{tealgreen}
From Nuclei to Patients: A Systematic Evaluation of Aggregation
Strategies for Bladder Cancer Recurrence Prediction}
\end{center}
\noindent\rule{\textwidth}{0.6pt}
\vspace{0.5em}

\section{Supplementary Methods}

\subsection{Neural Network Weight Budget}

To satisfy \texttt{nnet}'s weight limit, hidden units are capped
at $H_{\max} = \lfloor(3000-1)/(p+2)\rfloor$ for input
dimension $p$. For Representation~A ($p=460$), $H_{\max}=6$;
for Representation~B ($p=92$), $H_{\max}=10$. Five random
restarts per $(H,\lambda)$; best solution retained by training
loss.

\subsection{Handling Near-Zero-Variance Features in LDA}

Features f13--f23 and f27--f44 are identically zero for the
majority of nuclei. Before LDA, columns with within-class
SD~$<10^{-8}$ are removed. If fewer than two features remain,
the majority-class prediction is returned. All PCA+LDA inner-CV
steps use \texttt{tryCatch} to handle remaining singularities.

\section{Supplementary Figures}

\refstepcounter{suppfig}
\begin{figure}[H]
  \centering
  \includegraphics[width=\textwidth]{Fig4_FeatureImportance.png}
  \caption{\textbf{Figure S\arabic{suppfig}. Top 20 karyometric
    features by permutation importance} (Random Forest on
    Representation~A, all 84 patients, 500 trees). Colour by
    summary statistic: \textcolor{repA}{\textbf{blue}} = mean,
    \textcolor{repC}{\textbf{red}} = SD, orange = Q50 (median),
    purple = Q90. Features f54 and f53 dominate; the prevalence
    of SD and Q90 summaries confirms that within-patient nuclear
    variability is more predictive than mean morphology.}
  \label{fig:importance}
\end{figure}

\clearpage

\refstepcounter{suppfig}
\begin{figure}[H]
  \centering
  \includegraphics[width=\textwidth]{Fig5_Leakage.png}
  \caption{\textbf{Figure S\arabic{suppfig}. Leakage
    demonstration.} \textit{Panel (a):} M2 under valid
    patient-level LOOCV (MCE~$= 0.607$, blue) vs.\ leaky
    nucleus-level 10-fold CV (MCE~$= 0.268$, red). Leakage
    bias~$= 0.339$ (56\% overoptimism). Dashed line:
    majority-class baseline.
    \textit{Panel (b):} Nucleus-level recurrence risk scores
    by patient group. NR nuclei score near 0 (blue); R nuclei
    score near 0.8--0.9 (red). The strong nucleus-level separation
    explains the leakage: the scorer memorizes patient patterns,
    making held-out nuclei from the same patient trivially
    classifiable.}
  \label{fig:leakage}
\end{figure}

\clearpage

\refstepcounter{suppfig}
\begin{figure}[H]
  \centering
  \includegraphics[width=0.85\textwidth]{Fig7_Heatmap_MCE.png}
  \caption{\textbf{Figure S\arabic{suppfig}. MCE heatmap:
    Representation $\times$ Classifier} for M1--M12.
    Green = low MCE (better); red = high MCE. Representation~C
    row (top): uniformly red (MCE~$= 0.714$). Representation~B
    row (middle): consistently greener than Representation~A
    (bottom); best cell is M5, Lasso + Rep~B (MCE~$= 0.429$).
    Variation across rows (representations) exceeds variation
    across columns (classifiers), confirming aggregation strategy
    as the primary performance driver.}
  \label{fig:heatmap}
\end{figure}

\end{document}